\section{Preliminaries}\label{sec:preliminaries}

All vector norms, inner products, induced operator norms, and Frobenius norms are Euclidean. The symbol
$\lambda_k$ denotes $k$-dimensional Lebesgue measure in the original coefficient coordinates (and in the
nonpivot coordinates when the dimension is $N-1$); $\mathcal H^{N-1}$ denotes Euclidean Hausdorff measure on
sections. We use $\lambda_0([-R,R]^0)=1$ and regard $[-R,R]^0$ as the singleton containing the empty tuple.

For a fixed deterministic instance, $\mu$ and then an interval are selected only after all presentation data
$(\Theta,T,q,M,\Delta,N,R,\kappa,A,m,B)$ have been fixed. The capacities in the setup therefore always take
the interval supremum for each fixed law before the outer law supremum. Every probability is ordinary
probability for that fixed law; no independence is implicit.

For an affine section, $H_\theta$ is the actual set defined in the setup. When $F_0\equiv0$, it is the central
section $F(\theta)^\perp$, and $B_F$ denotes the principal block of $B$ indexed by $1,\ldots,N$. In a monic
specialization with degree $d$, the notation $p_\alpha$ always means the polynomial with the deterministic
leading coefficient stated in the theorem, while $\alpha$ contains only the lower coefficients.

The pointwise chart $T_j$ and insertion map $\Psi_j$ are the setting-defined original-coordinate maps. Empty
cells are allowed, and no lower bound on a pivot is part of the notation. These conventions are used in the
measurable chart statement and in the exact monic recovery.
